% ============================================================================
% PLANTILLA: Ejemplos de Tablas LaTeX para Diccionario de Datos
% ============================================================================
% Este archivo es SOLO REFERENCIA
% Copia estos ejemplos en documentacion/main.tex Sección 3 (Diccionario)
% ============================================================================

% ============================================================================
% EJEMPLO 1: Tabla Simple de Diccionario de Datos
% ============================================================================

% Para agregar una nueva tabla, copia esta estructura:

\subsection{Tabla: nombre_tabla}

\begin{center}
\begin{tabular}{|l|l|l|l|p{4cm}|}
    \hline
    \textbf{Campo} & \textbf{Tipo de Dato} & \textbf{NULL} & \textbf{Clave} & \textbf{Descripción} \\
    \hline
    campo_id & INT & NO & PK & Identificador único de la tabla \\
    \hline
    nombre & VARCHAR(100) & NO & & Nombre del registro \\
    \hline
    email & VARCHAR(100) & SÍ & UNIQUE & Correo electrónico único \\
    \hline
    fecha_creacion & TIMESTAMP & NO & & Fecha de creación del registro \\
    \hline
    activo & BOOLEAN & NO & & Estado del registro (Sí/No) \\
    \hline
\end{tabular}
\end{center}

% ============================================================================
% EJEMPLO 2: Tabla con Relaciones (Clave Foránea)
% ============================================================================

\subsection{Tabla: detalles_pedido}

\begin{center}
\begin{tabular}{|l|l|l|l|p{4cm}|}
    \hline
    \textbf{Campo} & \textbf{Tipo de Dato} & \textbf{NULL} & \textbf{Clave} & \textbf{Descripción} \\
    \hline
    detalles_pedido_id & INT & NO & PK & Identificador único \\
    \hline
    pedido_id & INT & NO & FK & Referencia a tabla pedidos \\
    \hline
    producto_id & INT & NO & FK & Referencia a tabla productos \\
    \hline
    cantidad & INT & NO & & Número de unidades compradas \\
    \hline
    precio_unitario & DECIMAL(10,2) & NO & & Precio por unidad al momento de compra \\
    \hline
    subtotal & DECIMAL(12,2) & NO & & cantidad × precio\_unitario \\
    \hline
\end{tabular}
\end{center}

\subsubsection{Relaciones}
\begin{itemize}
    \item \textbf{pedido\_id} es una Clave Foránea de \texttt{pedidos.pedido\_id}
    \item \textbf{producto\_id} es una Clave Foránea de \texttt{productos.producto\_id}
    \item Esta tabla implementa una relación N:M entre pedidos y productos
\end{itemize}

% ============================================================================
% EJEMPLO 3: Tabla con Múltiples Restricciones
% ============================================================================

\subsection{Tabla: inventario}

\begin{center}
\begin{tabular}{|l|l|l|l|p{4cm}|}
    \hline
    \textbf{Campo} & \textbf{Tipo de Dato} & \textbf{NULL} & \textbf{Clave} & \textbf{Descripción} \\
    \hline
    inventario_id & INT & NO & PK & Identificador único \\
    \hline
    producto_id & INT & NO & FK & Referencia a productos \\
    \hline
    cantidad & INT & NO & & Stock actual (CHECK: >= 0) \\
    \hline
    cantidad_minima & INT & NO & & Mínimo permitido antes de alerta \\
    \hline
    cantidad_maxima & INT & NO & & Máximo permitido en almacén \\
    \hline
    ubicacion & VARCHAR(50) & SÍ & & Ubicación física en almacén \\
    \hline
    fecha_actualizacion & TIMESTAMP & NO & & Última actualización de stock \\
    \hline
\end{tabular}
\end{center}

\subsubsection{Restricciones}
\begin{itemize}
    \item \textbf{cantidad}: CHECK (cantidad >= 0)
    \item \textbf{cantidad\_minima}: CHECK (cantidad\_minima < cantidad\_maxima)
    \item \textbf{producto\_id}: FOREIGN KEY UNIQUE (un producto por registro)
\end{itemize}

% ============================================================================
% EJEMPLO 4: Tabla de Auditoría
% ============================================================================

\subsection{Tabla: auditoria}

\begin{center}
\begin{tabular}{|l|l|l|l|p{4cm}|}
    \hline
    \textbf{Campo} & \textbf{Tipo de Dato} & \textbf{NULL} & \textbf{Clave} & \textbf{Descripción} \\
    \hline
    auditoria_id & INT & NO & PK & Identificador único de registro \\
    \hline
    tabla & VARCHAR(50) & NO & & Nombre de la tabla modificada \\
    \hline
    registro_id & INT & NO & & ID del registro modificado \\
    \hline
    operacion & ENUM(...) & NO & & INSERT, UPDATE o DELETE \\
    \hline
    datos_anteriores & JSON & SÍ & & Estado previo (para UPDATE) \\
    \hline
    datos_nuevos & JSON & SÍ & & Estado nuevo (para INSERT/UPDATE) \\
    \hline
    usuario & VARCHAR(100) & NO & & Quién realizó el cambio \\
    \hline
    fecha_cambio & TIMESTAMP & NO & & Cuándo se realizó el cambio \\
    \hline
\end{tabular}
\end{center}

% ============================================================================
% EJEMPLO 5: Tabla de Configuración/Parámetros
% ============================================================================

\subsection{Tabla: configuracion}

\begin{center}
\begin{tabular}{|l|l|l|l|p{4cm}|}
    \hline
    \textbf{Campo} & \textbf{Tipo de Dato} & \textbf{NULL} & \textbf{Clave} & \textbf{Descripción} \\
    \hline
    config_id & INT & NO & PK & Identificador único \\
    \hline
    parametro & VARCHAR(100) & NO & UNIQUE & Nombre del parámetro \\
    \hline
    valor & VARCHAR(255) & NO & & Valor del parámetro \\
    \hline
    tipo & ENUM('string', 'numero', 'booleano') & NO & & Tipo de dato esperado \\
    \hline
    descripcion & TEXT & SÍ & & Explicación del parámetro \\
    \hline
\end{tabular}
\end{center}

\subsubsection{Ejemplos de Parámetros}
\begin{itemize}
    \item \texttt{MAX\_ITEMS\_PER\_ORDER}: 100
    \item \texttt{IMPUESTO}: 12
    \item \texttt{PERMITE\_DEVOLUCIONES}: true
\end{itemize}

% ============================================================================
% EJEMPLO 6: Tabla de Relación N:M (Intermedia)
% ============================================================================

\subsection{Tabla: usuario\_rol}

\begin{center}
\begin{tabular}{|l|l|l|l|p{4cm}|}
    \hline
    \textbf{Campo} & \textbf{Tipo de Dato} & \textbf{NULL} & \textbf{Clave} & \textbf{Descripción} \\
    \hline
    usuario_id & INT & NO & PK/FK & Referencia a usuarios \\
    \hline
    rol_id & INT & NO & PK/FK & Referencia a roles \\
    \hline
    fecha_asignacion & TIMESTAMP & NO & & Cuándo se asignó el rol \\
    \hline
\end{tabular}
\end{center}

\subsubsection{Estructura Relacional}
\begin{itemize}
    \item \textbf{Clave Primaria Compuesta}: (usuario\_id, rol\_id)
    \item \textbf{usuario\_id}: Foreign Key → \texttt{usuarios.usuario\_id}
    \item \textbf{rol\_id}: Foreign Key → \texttt{roles.rol\_id}
    \item Implementa relación N:M: Un usuario puede tener múltiples roles
\end{itemize}

% ============================================================================
% GUÍA: Cómo Copiar Estas Tablas
% ============================================================================

% PASOS:
% 1. Abre documentacion/main.tex
% 2. Ve a la Sección 3 "Diccionario de Datos"
% 3. Busca el comentario "% TODO: [RESPONSABLE...]"
% 4. Copia UNA de estas tablas como referencia
% 5. Modifica:
%    - Nombre de la tabla (\subsection{Tabla: ...})
%    - Campos (nombres, tipos, restricciones)
%    - Descripción (explicar qué es cada campo)
% 6. Repite para cada tabla del MER
% 7. Recompila: pdflatex -interaction=nonstopmode main.tex

% ============================================================================
% TIPOS DE DATOS COMUNES EN MYSQL
% ============================================================================

% INT              - Número entero (4 bytes): -2,147,483,648 a 2,147,483,647
% BIGINT           - Número entero grande (8 bytes)
%
% DECIMAL(M,D)     - Decimal exacto: M=dígitos totales, D=decimales
%                     Ejemplo: DECIMAL(10,2) para precios
%
% VARCHAR(N)       - Texto variable: hasta N caracteres
% CHAR(N)          - Texto fijo: N caracteres
% TEXT             - Texto largo (mayor a 65KB)
%
% TIMESTAMP        - Fecha y hora: 'YYYY-MM-DD HH:MM:SS'
% DATE             - Solo fecha: 'YYYY-MM-DD'
% TIME             - Solo hora: 'HH:MM:SS'
%
% BOOLEAN          - TRUE / FALSE (en MySQL es equivalente a TINYINT(1))
%
% ENUM(...)        - Selección de valores: ENUM('opción1', 'opción2')
%
% JSON             - Datos JSON estructurados

% ============================================================================
% RESTRICCIONES COMUNES
% ============================================================================

% PRIMARY KEY (PK)        - Identifica unícamente cada fila
% FOREIGN KEY (FK)        - Referencia a otra tabla
% UNIQUE                  - Valor único en la columna
% NOT NULL                - Campo obligatorio
% DEFAULT valor           - Valor por defecto si no se especifica
% AUTO_INCREMENT          - Incrementa automáticamente
% CHECK (condición)       - Valida una condición
% UNSIGNED                - Solo números positivos

% ============================================================================
% FIN DE LA PLANTILLA
% ============================================================================
